% GENERATED FILE. Edit canonical lessons, then run npm run generate:books.
% canonical-source-hash: fnv1a64:b9917df1097a3e1a
% canonical-lessons: LA-C126-pallium, LA-C126-zona, LA-C126-pomum, LA-C126-cepa, LA-C126-alium

\chapter{Cloak, Belt, Fruit, Onion, Garlic}
\label{ch:la-cloak-belt-fruit-onion-garlic}
\hlchaptermodality{126}

\begin{chapteropening}
\textbf{By the end of this chapter:} \emph{I can say a cloak, a belt, a fruit, an onion and garlic.}

Complete the last lesson of chapter 126: I can say a cloak, a belt, a fruit, an onion and garlic.
\end{chapteropening}

\section[pallium, palliī]{pallium, palliī — a cloak}

\label{lesson:LA-C126-pallium}



\begin{quote}
Before the new one: say the Latin for a wagon, then the Latin for a long dress.
\end{quote}

\subsection*{You'll want to know: pallium, palliī}
\textbf{pallium, palliī} — \textquotedblleft{}a cloak\textquotedblright{}.

\begin{cousinweb}
\textbf{pallium} belongs with \emph{palla}, \textquotedblleft{}a woman's mantle\textquotedblright{} (Lewis \& Short); where the pair comes from is uncertain.
\end{cousinweb}

\begin{grammarlens}[title={What you've built}]
One more word for this chapter.
\end{grammarlens}

\subsection*{Your turn}
\begin{itemize}
\raggedright
  \item \emph{Say it:} \emph{pallium}
  \item \emph{Say it:} \emph{pallium}, once more
  \item \emph{Say it:} say \emph{pallium}
  \item \emph{Recall:} say the Latin for a wagon, then the Latin for a long dress, then say \emph{pallium} again
\end{itemize}

\begin{tcolorbox}[breakable,colback=teal!4,colframe=teal!35!black,title={Before you move on}]
What does pallium, palliī mean? (\textquotedblleft{}A cloak\textquotedblright{}.) Say it once more.
\end{tcolorbox}
\section[zōna, zōnae]{zōna, zōnae — a belt}

\label{lesson:LA-C126-zona}



\begin{quote}
Before the new one: say the Latin for a long dress, then the Latin for a cloak.
\end{quote}

\subsection*{You'll want to know: zōna, zōnae}
\textbf{zōna, zōnae} — \textquotedblleft{}a belt\textquotedblright{}.

\begin{cousinweb}
\textbf{zōna} was borrowed from Greek \emph{zōnē}, \textquotedblleft{}a belt, a girdle\textquotedblright{} (Lewis \& Short). English \emph{zone} is the same word: the climate belts of the earth.
\end{cousinweb}

\begin{grammarlens}[title={What you've built}]
One more word for this chapter.
\end{grammarlens}

\subsection*{Your turn}
\begin{itemize}
\raggedright
  \item \emph{Say it:} \emph{zōna}
  \item \emph{Say it:} \emph{zōna}, once more
  \item \emph{Say it:} say \emph{zōna}
  \item \emph{Recall:} say the Latin for a long dress, then the Latin for a cloak, then say \emph{zōna} again
\end{itemize}

\begin{tcolorbox}[breakable,colback=teal!4,colframe=teal!35!black,title={Before you move on}]
What does zōna, zōnae mean? (\textquotedblleft{}A belt\textquotedblright{}.) Say it once more.
\end{tcolorbox}
\section[pōmum, pōmī]{pōmum, pōmī — a fruit, an apple}

\label{lesson:LA-C126-pomum}



\begin{quote}
Before the new one: say the Latin for a cloak, then the Latin for a belt.
\end{quote}

\subsection*{You'll want to know: pōmum, pōmī}
\textbf{pōmum, pōmī} — \textquotedblleft{}a fruit, an apple\textquotedblright{}.

\begin{cousinweb}
\textbf{pōmum} is usually traced to Proto-Indo-European \emph{*h2epo}, \textquotedblleft{}off, away\textquotedblright{}, and \emph{*h1em-}, \textquotedblleft{}to take\textquotedblright{} (EtymDB 2.1), through an older \emph{*po-emom}, \textquotedblleft{}what is taken off, picked\textquotedblright{}. English \emph{pomegranate} is built on it.
\end{cousinweb}

\begin{grammarlens}[title={What you've built}]
One more word for this chapter.
\end{grammarlens}

\subsection*{Your turn}
\begin{itemize}
\raggedright
  \item \emph{Say it:} \emph{pōmum}
  \item \emph{Say it:} \emph{pōmum}, once more
  \item \emph{Say it:} say \emph{pōmum}
  \item \emph{Recall:} say the Latin for a cloak, then the Latin for a belt, then say \emph{pōmum} again
\end{itemize}

\begin{tcolorbox}[breakable,colback=teal!4,colframe=teal!35!black,title={Before you move on}]
What does pōmum, pōmī mean? (\textquotedblleft{}A fruit, an apple\textquotedblright{}.) Say it once more.
\end{tcolorbox}
\section[cēpa, cēpae]{cēpa, cēpae — an onion}

\label{lesson:LA-C126-cepa}



\begin{quote}
Before the new one: say the Latin for a belt, then the Latin for a fruit.
\end{quote}

\subsection*{You'll want to know: cēpa, cēpae}
\textbf{cēpa, cēpae} — \textquotedblleft{}an onion\textquotedblright{}.

\begin{cousinweb}
\textbf{cēpa}, also spelled \emph{caepa} (Lewis \& Short), has no sure origin. English \emph{chive} comes from it through French.
\end{cousinweb}

\begin{grammarlens}[title={What you've built}]
One more word for this chapter.
\end{grammarlens}

\subsection*{Your turn}
\begin{itemize}
\raggedright
  \item \emph{Say it:} \emph{cēpa}
  \item \emph{Say it:} \emph{cēpa}, once more
  \item \emph{Say it:} say \emph{cēpa}
  \item \emph{Recall:} say the Latin for a belt, then the Latin for a fruit, then say \emph{cēpa} again
\end{itemize}

\begin{tcolorbox}[breakable,colback=teal!4,colframe=teal!35!black,title={Before you move on}]
What does cēpa, cēpae mean? (\textquotedblleft{}An onion\textquotedblright{}.) Say it once more.
\end{tcolorbox}
\section[ālium, āliī]{ālium, āliī — garlic}

\label{lesson:LA-C126-alium}



\begin{quote}
Before the new one: say the Latin for a fruit, then the Latin for an onion.
\end{quote}

\subsection*{You'll want to know: ālium, āliī}
\textbf{ālium, āliī} — \textquotedblleft{}garlic\textquotedblright{}.

\begin{cousinweb}
\textbf{ālium}, also spelled \emph{allium} (Lewis \& Short), has no sure origin.
\end{cousinweb}

\begin{grammarlens}[title={What you've built}]
One more word for this chapter.
\end{grammarlens}

\subsection*{Your turn}
\begin{itemize}
\raggedright
  \item \emph{Say it:} \emph{ālium}
  \item \emph{Say it:} \emph{ālium}, once more
  \item \emph{Say it:} say \emph{ālium}
  \item \emph{Recall:} say the Latin for a fruit, then the Latin for an onion, then say \emph{ālium} again
\end{itemize}

\begin{tcolorbox}[breakable,colback=teal!4,colframe=teal!35!black,title={Before you move on}]
What does ālium, āliī mean? (\textquotedblleft{}Garlic\textquotedblright{}.) Say it once more.
\end{tcolorbox}
